\documentclass{article}
\usepackage{graphicx} % Required for inserting images
\usepackage{amsmath, amssymb}
\usepackage{multirow, makecell}

\usepackage{textcomp} % for copyright/registered symbols

\usepackage{bm} % for bold math symbols

\usepackage[a4paper, margin=2.5cm]{geometry} % change paper size
\usepackage{subcaption}

\usepackage{booktabs}
\usepackage{placeins} % for \FloatBarrier

\usepackage{float} % better image postioning wrt to txt

\usepackage{todonotes} % for TODOs
\let\oldtodo\todo % make \todo default to use inline
\renewcommand{\todo}[1]{\oldtodo[inline]{\detokenize{#1}}}


\usepackage{siunitx} % for units
\sisetup{per-mode=symbol}

 % for better tables/vertical centering
% https://tex.stackexchange.com/questions/7208/how-to-vertically-center-the-text-of-the-cells/611601#611601
\usepackage{tabularray}

\usepackage{textcomp} % for textrightarrow

\usepackage{xcolor} % for code blocks with color

% Code formatting
\definecolor{light-gray}{gray}{0.95}
\newcommand{\code}[1]{\colorbox{light-gray}{\texttt{\detokenize{#1}}}}

\usepackage{appendix}

% Enable links for figure numbers
\usepackage{hyperref} % note - hyperref likes to be last package
\hypersetup{
    colorlinks=true,
    linkcolor=blue,
}

\title{Research project}

\author{Jesse Li (6399053)}
\date{2026}


\begin{document}

\maketitle

\section{Background and Motivation}

\[\kappa = \frac{\sum_{j \in S_e} \rho_j w_j \lambda_j }{\sum_{j \in S_e} w_j / 2}\]

\section{Gradient floor and Hessian bound}

Together, a floor on $|\nabla t|$ and a bound on the Hessian of $t$ leave no interior critical point of $t$ inside any solid element.

\subsection*{1. Local quadratic model}
From central differences, estimate gradient $\mathbf{g}$ and Hessian $\mathbf{H}$:
\[
    \mathbf{g} = \begin{pmatrix} t_x \\ t_y \end{pmatrix}, \qquad
    \mathbf{H} = \begin{pmatrix} t_{xx} & t_{xy} \\ t_{xy} & t_{yy} \end{pmatrix}
\]
They define a quadratic model around the element centre $\mathbf{c}$:
\[
    t(\mathbf{c} + \mathbf{d}) \approx t(\mathbf{c}) + \mathbf{g}^\top \mathbf{d} + \tfrac{1}{2} \mathbf{d}^\top \mathbf{H} \mathbf{d}
    \quad \Longrightarrow \quad
    \nabla t(\mathbf{c} + \mathbf{d}) = \mathbf{g} + \mathbf{H} \mathbf{d}
\]

\subsection*{2. The gradient changes little across the element}
Constraint: Critical point of this model should be outside the element, so:
\[\left| \nabla t(\mathbf{c} + \mathbf{d}) \right| = \left|\mathbf{g} + \mathbf{H} \mathbf{d}\right|> 0\]

From Pythagoras, $|\mathbf{d}| \le h/\sqrt{2}$. Then:
\[
    |\mathbf{H} \mathbf{d}|
    \le \|\mathbf{H}\|_2 \, |\mathbf{d}|
    \le \|\mathbf{H}\|_F \, \frac{h}{\sqrt{2}},
    \qquad
    \|\mathbf{H}\|_F = \sqrt{t_{xx}^2 + t_{yy}^2 + 2 t_{xy}^2}
\]
So, constrain
\[
    |\mathbf{g}| > \|\mathbf{H}\|_F \, \frac{h}{\sqrt{2}}
\]

\subsection*{4. The two constraints}
Let $m$ be the median of $|\mathbf{g}|$ over the mesh, held constant in the sensitivities. For a solid element, with floor fraction $f$ and smoothness length $L$:
\[
    \text{floor:} \quad |\mathbf{g}| \ge f\, m,
    \qquad
    \text{smoothness:} \quad \|\mathbf{H}\|_F \le \frac{m}{L}
\]

\subsection*{5. Combined}
A critical point would need
\[
    f\, m \le |\mathbf{g}| \le \|\mathbf{H}\|_F \, \frac{h}{\sqrt{2}} \le \frac{m}{L} \, \frac{h}{\sqrt{2}}
\]
so none exists in the element if
\[
    f > \frac{h}{\sqrt{2}\, L}
    \quad \Longleftrightarrow \quad
    L > \frac{h}{\sqrt{2}\, f}
\]

\section{Heat equation}
$$\nabla\cdot\big(k(x)\,\nabla u\big) = 0, \qquad 0 < \lambda \le k(x) \le \Lambda .$$

Let $k(x) = e^{w(x)}$.
\[0 = \nabla \cdot (k \nabla u) = k \Delta u + \nabla k \cdot \nabla u = e^w \Delta u + e^w \nabla w \cdot \nabla u\]
\[\nabla u \cdot \nabla w = -\Delta u\]

\end{document}


\end{document}
